\section{Fold content and the radial graph}\label{sec:fold}

The trapping annulus of \cref{thm:annulus} confines $\partial\Etil$ to a
thin shell $\{\rho_i\le\abs{x-z}\le\rho_e\}$ of width $\eta=\rho_e-\rho_i$,
so the inset has \emph{small amplitude}. To convert this into a radial
\emph{graph} about the centre $z$ we must rule out radial folds: rays from
$z$ meeting $\partial\Etil$ more than once. This section bounds the total
\emph{fold content} --- the obstruction to being a graph --- by the
geometric deficit, with \emph{no} square-root loss from the quantitative
isoperimetric inequality, and then realises $\Etil$ as a measurable radial
graph up to a controlled symmetric difference.

Throughout, $\Etil$ is the hole-filled inset of \cref{def:inset,lem:fill}
(\cref{conv:fill}), connected and globally hole-free, so the global
exterior-connectedness needed to invoke the trapping annulus of
\cref{thm:annulus} holds unconditionally (\cref{lem:fill}\ref{it:fill-top}).
We write
with $\mu_E:=\abs\Etil$, $\hat r:=\sqrt{\mu_E/\pi}$, isoperimetric deficit
$\diso:=\diso(\Etil)$, perimeter excess $\Exc:=\Per(\Etil)-2\pi\hat r
=2\pi\hat r\,\diso$, and (\cref{thm:annulus}) trapping
\begin{equation}\label{eq:fold-annulus}
   B_{\rho_i}(z)\subset\Etil\subset B_{\rho_e}(z),\qquad
   \rho_i\le\hat r\le\rho_e,\qquad
   \eta:=\rho_e-\rho_i\le C\,(\Exc+a)^{1/2},
\end{equation}
where $z$ is the optimal centre and $a:=\abs{\Etil\triangle B_{\hat r}(z)}$.
In particular $z$ lies in the interior of $\Etil$ (indeed
$B_{\rho_i}(z)\subset\Etil$), and $\partial\Etil$ avoids the singular point:
\begin{equation}\label{eq:fold-away}
   \partial\Etil\subset\{x:\rho_i\le\abs{x-z}\le\rho_e\}.
\end{equation}
We write $e_r:=(x-z)/\abs{x-z}$, $e_\theta:=e_r^\perp$ for the unit radial
and angular vectors at $x\ne z$, $\nu$ for the outward unit normal on the
reduced boundary $\partial^*\Etil$, and
\[
   \nu_r:=\nu\cdot e_r
\]
for the radial part of $\nu$. (As is standard for sets of finite perimeter,
all boundary integrals are over the reduced boundary $\partial^*\Etil$, on
which $\nu$ is defined $\HH^1$-a.e.; we write $\partial\Etil$ for it
throughout, noting that for the hole-free finite-perimeter set $\Etil$ the
reduced and essential boundaries coincide $\HH^1$-a.e.\ by the De
Giorgi--Federer structure theorem.) For $\theta\in S^1$ let
\begin{equation}\label{eq:fold-M}
   M(\theta):=\HH^0\bigl(\partial\Etil\cap R_\theta\bigr),\qquad
   R_\theta:=\{z+\rho e^{i\theta}:\rho>0\}
\end{equation}
be the number of crossings of $\partial\Etil$ by the open ray in direction
$\theta$. By the one-dimensional structure of slices of finite-perimeter
sets \cite[\S2.4, \S3.11]{AFP2000}, for a.e.\ $\theta$ the radial slice
\[
   \Etil_\theta:=\{\rho>0:z+\rho e^{i\theta}\in\Etil\}
\]
is, up to an $\HH^1$-null set, a finite union of intervals; since
$B_{\rho_i}(z)\subset\Etil\subset B_{\rho_e}(z)$, the innermost interval
starts at $0$ and all endpoints lie in $[\rho_i,\rho_e]$, so
\begin{equation}\label{eq:fold-slice}
   \Etil_\theta=(0,r_1)\cup(r_2,r_3)\cup\cdots\cup(r_{2m},r_{2m+1}),
   \quad \rho_i\le r_1\le\cdots\le r_{2m+1}\le\rho_e,
   \quad M(\theta)=2m(\theta)+1.
\end{equation}
The intervals $(r_{2i},r_{2i+1})$, $i\ge1$, are the \emph{caps};
$m(\theta)$ is the number of caps on the ray, and $\Etil$ is a radial graph
about $z$ exactly when $M\equiv1$ (no caps).

\subsection{Point-source flux identities}

\begin{lemma}[Two-flux identities]\label{lem:two-flux}
Let $E\subset\R^2$ be a set of finite perimeter and suppose that, for some
$r_0>0$,
\[
   B_{r_0}(z)\subset E,\qquad
   \partial^* E\subset\{x:\abs{x-z}\ge r_0\}.
\]
With $e_r=(x-z)/\abs{x-z}$, $\nu_r=\nu\cdot e_r$, and
$M(\theta)$ as in \eqref{eq:fold-M}, the following hold:
\begin{subequations}\label{eq:two-flux}
\begin{align}
   \textup{(i)}\quad
   &\int_{\partial E}\frac{\abs{\nu_r}}{\abs{x-z}}\,d\HH^1
   =\int_{S^1}M(\theta)\,d\theta,
   \label{eq:two-flux-i}\\[2pt]
   \textup{(ii)}\quad
   &\int_{\partial E}\frac{\nu_r}{\abs{x-z}}\,d\HH^1=2\pi,
   \label{eq:two-flux-ii}\\[2pt]
   \textup{(iii)}\quad
   &\int_{\partial E}\nu_r\,d\HH^1
   =\int_E\frac{dx}{\abs{x-z}}=:I_E .
   \label{eq:two-flux-iii}
\end{align}
\end{subequations}
We refer to the three identities collectively as \eqref{eq:two-flux}.
\end{lemma}

\begin{proof}
We may translate so that $z=0$; then $e_r=x/\abs x$ and $\abs{x-z}=\abs x$.
Set $A:=\{x:r_0\le\abs x\}$, a closed set containing $\partial^*E$ on which
all three integrands are smooth and bounded. Since $B_{r_0}(0)\subset E$,
for every $0<\eps<r_0$ the circle $\partial B_\eps(0)$ lies in the
measure-theoretic interior $E^{(1)}$.

\emph{(i) The angular coarea identity.}
Let $\Theta(x):=x/\abs x\in S^1$ be the angular projection.  On the annulus
$A$ this map is Lipschitz.  The reduced boundary $\partial^*E$ is a
countably $1$-rectifiable set contained in $A$; let $\tau$ denote a unit
tangent to $\partial^*E$ (defined $\HH^1$-a.e., orthogonal to $\nu$). The
tangential Jacobian of $\Theta$ along $\partial^*E$ is
\[
   J_{\partial^*E}\Theta
   =\frac{\abs{e_\theta\cdot\tau}}{\abs x}.
\]
Since $\{e_r,e_\theta\}$ and $\{\nu,\tau\}$ are each orthonormal bases of
$\R^2$ with $e_r\perp e_\theta$ and $\nu\perp\tau$, the angle between
$e_\theta$ and $\tau$ equals the angle between $e_r$ and $\nu$; hence
$\abs{e_\theta\cdot\tau}=\abs{e_r\cdot\nu}=\abs{\nu_r}$, so
\[
   J_{\partial^*E}\Theta=\frac{\abs{\nu_r}}{\abs x}.
\]
The coarea formula for the Lipschitz map
$\Theta\colon\partial^*E\to S^1$ on the $1$-rectifiable set
$\partial^*E$ \cite[Theorem~2.93]{AFP2000} gives
\[
   \int_{\partial^*E}\frac{\abs{\nu_r}}{\abs x}\,d\HH^1
   =\int_{\partial^*E}J_{\partial^*E}\Theta\,d\HH^1
   =\int_{S^1}\HH^0\bigl(\partial^*E\cap \Theta^{-1}(\theta)\bigr)\,d\theta
   =\int_{S^1}M(\theta)\,d\theta,
\]
because $\Theta^{-1}(\theta)\cap A=R_\theta\cap A\supset
\partial^*E\cap R_\theta$ and $\partial^*E\subset A$. This is
\eqref{eq:two-flux-i}.

\emph{(ii) The point-source field, singularity excised.}
Consider the vector field $X(x):=x/\abs x^2$, defined and smooth on
$\R^2\setminus\{0\}$, with
\[
   \divg X=\frac{2}{\abs x^2}-\frac{2}{\abs x^2}=0
   \qquad(x\ne0)
\]
(directly: $\partial_i(x_i\abs x^{-2})=\abs x^{-2}+x_i(-2x_i\abs x^{-4})$,
summing $i=1,2$ gives $2\abs x^{-2}-2\abs x^{-2}=0$). On $\partial^*E$ we
have $X\cdot\nu=(x\cdot\nu)/\abs x^2=\nu_r/\abs x$. Fix $0<\eps<r_0$ and
apply the Gauss--Green theorem for sets of finite perimeter
\cite[Theorem~5.16]{EvansGariepy2015} to the field $X$, which is $C^1$ and
bounded on the closed set $\{\abs x\ge\eps/2\}$, over the finite-perimeter
set $E_\eps:=E\setminus\overline{B_\eps(0)}$. Since $\partial B_\eps\subset
E^{(1)}$ (shown above), the reduced boundary $\partial^*E_\eps$ splits
($\HH^1$-a.e.) into $\partial^*E$ (with outward normal $\nu$, and on which
$X$ is integrable since $\partial^*E\subset A$) and the inner circle
$\partial B_\eps$ (with outward-for-$E_\eps$ normal $-x/\eps$):
\[
   0=\int_{E_\eps}\divg X\,dx
   =\int_{\partial^*E}X\cdot\nu\,d\HH^1
   +\int_{\partial B_\eps}X\cdot\Bigl(-\frac{x}{\eps}\Bigr)\,d\HH^1 .
\]
On $\partial B_\eps$, $X=x/\eps^2$ and $X\cdot(-x/\eps)=-\abs x^2/\eps^3
=-1/\eps$, so the second term equals $-\eps^{-1}\HH^1(\partial B_\eps)
=-\eps^{-1}\cdot2\pi\eps=-2\pi$, independent of $\eps$. Hence
\[
   \int_{\partial^*E}\frac{\nu_r}{\abs x}\,d\HH^1
   =\int_{\partial^*E}X\cdot\nu\,d\HH^1=2\pi,
\]
which is \eqref{eq:two-flux-ii}. (No limit in $\eps$ is even needed: the
left side does not depend on $\eps$, and the inner-circle flux is exactly
$2\pi$ for every admissible $\eps$.)

\emph{(iii) The radial field, singularity integrable.}
Consider $Y(x):=e_r=x/\abs x$, smooth on $\R^2\setminus\{0\}$, with
\[
   \divg Y=\sum_i\partial_i\frac{x_i}{\abs x}
   =\sum_i\Bigl(\frac1{\abs x}-\frac{x_i^2}{\abs x^3}\Bigr)
   =\frac2{\abs x}-\frac{\abs x^2}{\abs x^3}=\frac1{\abs x}
   \qquad(x\ne0).
\]
Again fix $0<\eps<r_0$ and apply Gauss--Green to the $C^1$ bounded field $Y$
over $E_\eps=E\setminus\overline{B_\eps(0)}$:
\[
   \int_{E_\eps}\frac{dx}{\abs x}
   =\int_{E_\eps}\divg Y\,dx
   =\int_{\partial^*E}Y\cdot\nu\,d\HH^1
   +\int_{\partial B_\eps}Y\cdot\Bigl(-\frac x\eps\Bigr)\,d\HH^1 .
\]
On $\partial B_\eps$, $Y=x/\eps$ and $Y\cdot(-x/\eps)=-\abs x^2/\eps^2=-1$,
so the inner term is $-\HH^1(\partial B_\eps)=-2\pi\eps$. Thus
\[
   \int_{\partial^*E}\nu_r\,d\HH^1
   =\int_{E_\eps}\frac{dx}{\abs x}+2\pi\eps .
\]
Now let $\eps\downarrow0$. The left side is fixed. On the right,
$\abs x^{-1}\in L^1_{\mathrm{loc}}(\R^2)$ (in polar coordinates
$\int_{B_R}\abs x^{-1}\,dx=2\pi R<\infty$) and $E$ has finite measure, so
$\abs x^{-1}\mathbf1_E\in L^1(\R^2)$; by dominated convergence
$\int_{E_\eps}\abs x^{-1}\,dx\to\int_E\abs x^{-1}\,dx=I_E$, while
$2\pi\eps\to0$. Therefore $\int_{\partial^*E}\nu_r\,d\HH^1=I_E$, which is
\eqref{eq:two-flux-iii}. The constant $I_E$ is finite because
$\abs x^{-1}\mathbf1_E\in L^1$.
\end{proof}

\begin{remark}[On the singularity]\label{rem:singularity}
The two divergence-theorem identities are genuinely different at the source.
In \eqref{eq:two-flux-ii} the field $X=x/\abs x^2$ has a non-integrable
$\abs x^{-1}$ growth at $0$ and a non-zero residual flux $2\pi$ through any
small circle; excising $B_\eps$ is mandatory, and the $2\pi$ is precisely
that residue (the planar Green-function source). In \eqref{eq:two-flux-iii}
the field $Y=e_r$ is bounded, $\divg Y=\abs x^{-1}$ is integrable, and the
inner-circle flux $2\pi\eps\to0$ vanishes; the excision is a convenience and
the resulting $I_E$ is a finite, well-defined number. We invoke
\cref{lem:two-flux} for $E=\Etil$, whose reduced boundary satisfies the
hypothesis $\partial^*\Etil\subset\{\abs{x-z}\ge\rho_i\}$ with $r_0=\rho_i$
by \eqref{eq:fold-away}, and whose interior contains $z$ since
$B_{\rho_i}(z)\subset\Etil$.
\end{remark}

\subsection{Sharp fold content}

The fold content is the deviation of $M$ from the graph value $1$. We bound
its total mass with no quantitative-isoperimetric square-root loss; this is
the crux of the sharp reduction.

\begin{proposition}[Sharp fold content]\label{prop:fold}
With the inset $\Etil$ and the trapping data \eqref{eq:fold-annulus},
\begin{equation}\label{eq:foldsharp}
   \int_{S^1}\bigl(M(\theta)-1\bigr)\,d\theta
   \ =\ 2\!\int_{\partial\Etil\cap\{\nu_r<0\}}\!
        \frac{\abs{\nu_r}}{\abs{x-z}}\,d\HH^1
   \ \le\ \frac1{\rho_i}\Bigl(\Exc+\frac{a\,\eta}{\rho_i\hat r}\Bigr)
   \ \le\ C\bigl(\Exc+a\eta\bigr),
\end{equation}
with $C$ absolute. In particular the integrand is non-negative, so
$M\ge1$ a.e., and
\begin{equation}\label{eq:foldangles}
   \bigl|\{\theta\in S^1:M(\theta)\ge3\}\bigr|
   \ \le\ \tfrac12\!\int_{S^1}\bigl(M-1\bigr)\,d\theta
   \ \le\ C\bigl(\Exc+a\eta\bigr).
\end{equation}
\end{proposition}

\begin{proof}
\emph{The fold identity.}
Subtract \eqref{eq:two-flux-ii} from \eqref{eq:two-flux-i} of
\cref{lem:two-flux} (legitimate by \cref{rem:singularity}); since
$\int_{S^1}1\,d\theta=2\pi$ equals the right-hand side of
\eqref{eq:two-flux-ii},
\[
   \int_{S^1}\bigl(M-1\bigr)\,d\theta
   =\int_{\partial\Etil}\frac{\abs{\nu_r}-\nu_r}{\abs{x-z}}\,d\HH^1
   =2\!\int_{\partial\Etil\cap\{\nu_r<0\}}\!
       \frac{\abs{\nu_r}}{\abs{x-z}}\,d\HH^1\ \ge0,
\]
using $\abs{\nu_r}-\nu_r=0$ where $\nu_r\ge0$ and $=2\abs{\nu_r}$ where
$\nu_r<0$. This is the first equality of \eqref{eq:foldsharp}. (Pointwise
$M(\theta)=2m(\theta)+1\ge1$ a.e.\ by the slice structure
\eqref{eq:fold-slice}, since every slice contains $(0,r_1)$; the content of
the identity is the \emph{global} bound on the excess $\int(M-1)$, which is
exactly twice the inward-pointing radial flux.)

\emph{From the radial fold flux to $P_r-I_{\Etil}$.}
On the folded part $\partial\Etil\cap\{\nu_r<0\}$ we bound
$\abs{x-z}\ge\rho_i$ (by \eqref{eq:fold-away}), so
\begin{equation}\label{eq:fold-step1}
   \int_{S^1}\bigl(M-1\bigr)\,d\theta
   \ \le\ \frac2{\rho_i}\!\int_{\partial\Etil\cap\{\nu_r<0\}}\!
          \abs{\nu_r}\,d\HH^1
   \ =\ \frac{2F_r}{\rho_i},\qquad
   F_r:=\!\int_{\partial\Etil\cap\{\nu_r<0\}}\!\abs{\nu_r}\,d\HH^1 .
\end{equation}
Write $P_r:=\int_{\partial\Etil}\abs{\nu_r}\,d\HH^1$ for the total radial
perimeter. Splitting according to the sign of $\nu_r$,
\[
   \int_{\partial\Etil}\nu_r\,d\HH^1
   =\int_{\{\nu_r\ge0\}}\abs{\nu_r}\,d\HH^1
    -\int_{\{\nu_r<0\}}\abs{\nu_r}\,d\HH^1
   =\bigl(P_r-F_r\bigr)-F_r=P_r-2F_r .
\]
By \eqref{eq:two-flux-iii}, the left side is $I_{\Etil}$, hence
\begin{equation}\label{eq:fold-Fr}
   2F_r=P_r-I_{\Etil} .
\end{equation}

\emph{Upper bound on $P_r$.}
Since $\abs{\nu_r}=\abs{\nu\cdot e_r}\le\abs\nu=1$ pointwise, the radial
perimeter is dominated by the perimeter:
\[
   P_r=\int_{\partial\Etil}\abs{\nu_r}\,d\HH^1
   \le\int_{\partial\Etil}1\,d\HH^1=\Per(\Etil)=2\pi\hat r+\Exc .
\]

\emph{The cancellation: lower bound on $I_{\Etil}$.}
Here is the sharp step. Because $\abs\Etil=\mu_E=\pi\hat r^2
=\abs{B_{\hat r}(z)}$, the indicators have equal mass, so
$\int(\mathbf1_{\Etil}-\mathbf1_{B_{\hat r}(z)})\,dx=0$ and we may subtract
\emph{any} constant from the kernel without changing the integral. Subtract
the constant $1/\hat r$:
\[
   I_{\Etil}-\int_{B_{\hat r}(z)}\frac{dx}{\abs{x-z}}
   =\int\bigl(\mathbf1_{\Etil}-\mathbf1_{B_{\hat r}(z)}\bigr)
        \frac{dx}{\abs{x-z}}
   =\int\bigl(\mathbf1_{\Etil}-\mathbf1_{B_{\hat r}(z)}\bigr)
        \Bigl(\frac1{\abs{x-z}}-\frac1{\hat r}\Bigr)\,dx,
\]
the last equality because $\int(\mathbf1_{\Etil}-\mathbf1_{B_{\hat r}(z)})
\cdot\hat r^{-1}\,dx=\hat r^{-1}(\abs\Etil-\abs{B_{\hat r}(z)})=0$. In polar
coordinates,
$\int_{B_{\hat r}(z)}\abs{x-z}^{-1}\,dx
=\int_0^{\hat r}\!\!\int_{S^1}s^{-1}\,s\,d\theta\,ds=2\pi\hat r$, so
\begin{equation}\label{eq:fold-Idiff}
   I_{\Etil}-2\pi\hat r
   =\int_{\Etil\triangle B_{\hat r}(z)}
       \pm\Bigl(\frac1{\abs{x-z}}-\frac1{\hat r}\Bigr)\,dx,
\end{equation}
the integrand being supported on the symmetric difference
$\Etil\triangle B_{\hat r}(z)$ (with sign $+$ on $\Etil\setminus
B_{\hat r}(z)$ and $-$ on $B_{\hat r}(z)\setminus\Etil$). Now
$\Etil\triangle B_{\hat r}(z)$ lies in the closed annulus
$\{\rho_i\le\abs{x-z}\le\rho_e\}$: indeed
$B_{\rho_i}(z)\subset\Etil\cap B_{\hat r}(z)$ (as $\rho_i\le\hat r$), so
neither $\Etil$ nor $B_{\hat r}(z)$ differs from the other inside
$B_{\rho_i}(z)$; and $\Etil\cup B_{\hat r}(z)\subset B_{\rho_e}(z)$ (as
$\hat r\le\rho_e$ and $\Etil\subset B_{\rho_e}(z)$), so neither differs
outside $B_{\rho_e}(z)$. On that annulus, with both $\abs{x-z}$ and
$\hat r$ in $[\rho_i,\rho_e]$,
\begin{equation}\label{eq:fold-kernel}
   \Bigl|\frac1{\abs{x-z}}-\frac1{\hat r}\Bigr|
   =\frac{\bigl|\hat r-\abs{x-z}\bigr|}{\hat r\,\abs{x-z}}
   \le\frac{\rho_e-\rho_i}{\hat r\,\rho_i}=\frac{\eta}{\rho_i\hat r}.
\end{equation}
Therefore, taking absolute values in \eqref{eq:fold-Idiff} and using
\eqref{eq:fold-kernel},
\begin{equation}\label{eq:Icancel}
   \bigl|I_{\Etil}-2\pi\hat r\bigr|
   \ \le\ \frac{\eta}{\rho_i\hat r}\,
          \bigl|\Etil\triangle B_{\hat r}(z)\bigr|
   \ =\ \frac{a\,\eta}{\rho_i\hat r}.
\end{equation}

\emph{Conclusion.}
Combining \eqref{eq:fold-Fr} with the bound on $P_r$ and the lower bound
$I_{\Etil}\ge2\pi\hat r-a\eta/(\rho_i\hat r)$ from \eqref{eq:Icancel},
\[
   2F_r=P_r-I_{\Etil}
   \le\bigl(2\pi\hat r+\Exc\bigr)
      -\Bigl(2\pi\hat r-\frac{a\eta}{\rho_i\hat r}\Bigr)
   =\Exc+\frac{a\eta}{\rho_i\hat r}.
\]
Inserting into \eqref{eq:fold-step1} yields the middle inequality of
\eqref{eq:foldsharp}. Finally, in the optimised inset regime of
\cref{thm:reduction}, \eqref{eq:close-inset} gives
$\mu_E=\pi-O(\sqrt{\dT})$, hence $\hat r\ge\tfrac34$ after shrinking
$\delta_\star$, and \cref{thm:annulus} gives
$\eta\le C\dT^{1/8}\le\tfrac14$. Thus $\rho_i\ge\hat r-\eta\ge\tfrac12$
(using
$\rho_e=\rho_i+\eta$ and $\rho_i\le\hat r\le\rho_e$). Therefore
$1/\rho_i\le2$ and $1/(\rho_i^2\hat r)\le1/(\tfrac14\cdot\tfrac34)\le6$,
giving the absolute bound $C(\Exc+a\eta)$. The fold-angle estimate \eqref{eq:foldangles} follows
because $M-1\ge2$ on $\{M\ge3\}$ (as $M=2m+1$ is odd, $M\ge3$ means
$m\ge1$), so $2\abs{\{M\ge3\}}\le\int_{\{M\ge3\}}(M-1)\le\int_{S^1}(M-1)$.
\end{proof}

\begin{remark}[Why the factor $\eta$ matters: the improvement over the
naive bound]\label{rem:eta-matters}
The naive estimate pulls $\abs{x-z}^{-1}\le\rho_i^{-1}$ out of the integral
defining $I_{\Etil}-2\pi\hat r$ before subtracting any constant, giving
\[
   \bigl|I_{\Etil}-2\pi\hat r\bigr|
   =\Bigl|\int\bigl(\mathbf1_{\Etil}-\mathbf1_{B_{\hat r}(z)}\bigr)
        \tfrac1{\abs{x-z}}\Bigr|
   \le\frac1{\rho_i}\,\abs{\Etil\triangle B_{\hat r}(z)}
   =\frac{a}{\rho_i}\le Ca,
\]
hence $\int(M-1)\le C(\Exc+a)$. The two estimates differ by the factor
$\eta=\rho_e-\rho_i$ on the $a$-term. This is decisive. By the quantitative
isoperimetric inequality (\cref{thm:fmp}), $a\le C_{\mathrm F}\sqrt\diso\,
\mu_E$ carries an unavoidable \emph{square-root} of the isoperimetric
deficit; with $\diso\sim\dT/\tau\sim\sqrt\dT$ at the boundary scale this
gives only $a\sim\dT^{1/4}$, and the naive fold content
$\Exc+a\sim\dT^{1/4}$ would cap the whole argument at $\dT^{1/4}$ closeness.
The cancellation \eqref{eq:Icancel} replaces $a$ by $a\eta$; since
$\eta\le C(\Exc+a)^{1/2}\le Ca^{1/2}\sim\dT^{1/8}$, the product
$a\eta\sim\dT^{3/8}$, and through \cref{cor:gd} the graph defect becomes
$\gd\sim\eta\,a\eta\sim\dT^{1/8}\cdot\dT^{3/8}=\dT^{1/2}$, matching the
closeness $O(\sqrt\dT)$. The single factor $\eta$ --- the legitimacy of
subtracting the constant $1/\hat r$, because $\abs\Etil=\abs{B_{\hat r}(z)}$,
together with the confinement of $\Etil\triangle B_{\hat r}(z)$ to the thin
$\eta$-annulus --- is exactly what removes the FMP square-root from the fold
content and makes the reduction sharp.
\end{remark}

\subsection{The graph defect}

The \emph{graph defect} of $\Etil$ about $z$,
\[
   \gd(\Etil):=\inf_{\rho\colon S^1\to(0,\infty)}
      \bigl|\Etil\,\triangle\,\{z+r\omega:0\le r<\rho(\omega)\}\bigr|,
\]
measures how far $\Etil$ is, in $L^1$, from being a radial graph about $z$.
The fold content controls it directly.

\begin{corollary}[Graph defect]\label{cor:gd}
The graph defect of $\Etil$ about the centre $z$ satisfies
\begin{equation}\label{eq:gd}
   \gd(\Etil)\ \le\ \rho_e\,\eta\!\int_{S^1}\!m(\theta)\,d\theta
   \ =\ \tfrac12\,\rho_e\,\eta\!\int_{S^1}\!\bigl(M-1\bigr)\,d\theta
   \ \le\ C\,\eta\,\bigl(\Exc+a\eta\bigr),
\end{equation}
with $C$ absolute.
\end{corollary}

\begin{proof}
Use the \emph{core profile} $\rho_0(\theta):=r_1(\theta)\in[\rho_i,\rho_e]$,
the innermost crossing radius from \eqref{eq:fold-slice}; it is defined for
a.e.\ $\theta$. Set
$G_0:=\{z+r e^{i\theta}:0\le r<\rho_0(\theta)\}$. By construction the slice
of $G_0$ in direction $\theta$ is $(0,\rho_0(\theta))=(0,r_1(\theta))
\subset\Etil_\theta$, so $G_0\subset\Etil$ (up to $\HH^1$-null directions),
and the slice of $\Etil\setminus G_0$ is exactly the union of caps
$(r_2,r_3)\cup\cdots\cup(r_{2m},r_{2m+1})$. By the polar-coordinate area
formula (Fubini on $\R^2\setminus\{z\}$),
\[
   \bigl|\Etil\setminus G_0\bigr|
   =\int_{S^1}\sum_{i=1}^{m(\theta)}\int_{r_{2i}}^{r_{2i+1}}\!r\,dr\,d\theta .
\]
On each ray with $m(\theta)\ge1$, all cap endpoints lie in $[\rho_i,\rho_e]$,
so the caps have total radial length $\sum_i(r_{2i+1}-r_{2i})\le
r_{2m+1}-r_2\le\rho_e-\rho_i=\eta$, and each cap sits at radius
$r\le\rho_e$; hence
\[
   \sum_{i=1}^{m(\theta)}\int_{r_{2i}}^{r_{2i+1}}\!r\,dr
   \ \le\ \rho_e\sum_{i=1}^{m(\theta)}(r_{2i+1}-r_{2i})
   \ \le\ \rho_e\,\eta\,\mathbf1_{\{m(\theta)\ge1\}}.
\]
Therefore
\[
   \gd(\Etil)\le\bigl|\Etil\triangle G_0\bigr|=\bigl|\Etil\setminus G_0\bigr|
   \le\rho_e\,\eta\,\bigl|\{\theta:m(\theta)\ge1\}\bigr|
   \le\rho_e\,\eta\!\int_{S^1}m\,d\theta,
\]
where the last step uses $\mathbf1_{\{m\ge1\}}\le m$ (as $m$ is a
non-negative integer). Since $M=2m+1$, $\int_{S^1}m\,d\theta=\tfrac12
\int_{S^1}(M-1)\,d\theta$, which is the displayed equality; and
\cref{prop:fold} together with $\rho_e\le C\hat r\le C$ (for
$\dT\le\delta_\star$) gives the final bound
$\le C\,\eta(\Exc+a\eta)$.
\end{proof}

\subsection{Realization as a measurable radial graph}

The core profile $\rho_0=r_1$ is, a priori, only \emph{measurable} in
$\theta$ (it is defined ray by ray from the one-dimensional slices, with no
modulus of continuity and no bounded variation available --- a cap on a ray
inserts a jump that the slice structure does not control). We do not attempt
to regularise it. Instead we record that the measurable core graph $G_0$ is
\emph{already} an admissible input to the nearly-spherical estimate of
\cref{sec:nearly}, whose hypotheses --- a bounded open star-shaped set
trapped between two concentric circles, with measurable radial profile and
small $L^\infty$ amplitude --- ask for no slope, no continuity, and no
bounded variation. This removes the regularity step entirely.

\begin{proposition}[Measurable radial-graph realization]\label{prop:realization}
Let $\rho_0=r_1\colon S^1\to[\rho_i,\rho_e]$ be the innermost-crossing core
profile of \cref{cor:gd}, and let
\[
   G_0:=\{z+r\omega(\theta):\ \theta\in S^1,\ 0\le r<\rho_0(\theta)\}
\]
be the associated core graph, taken in its open star-shaped representative
\textup{(}equivalently, replace $\rho_0$ by its lower-semicontinuous
envelope along rays, which alters $G_0$ only on an $\HH^2$-null set and hence
changes none of the quantities below\textup{)}. Then:
\begin{enumerate}[label=\textup{(\roman*)},leftmargin=2.4em]
\item\label{it:real-meas} $\rho_0$ is measurable with values in
$[\rho_i,\rho_e]$, and $G_0$ is a bounded open set, star-shaped about $z$,
with
\begin{equation}\label{eq:real-trap}
   B_{\rho_i}(z)\subset G_0\subset B_{\rho_e}(z);
\end{equation}
\item\label{it:real-amp} writing $\hat r_{G_0}:=\sqrt{\abs{G_0}/\pi}\in
[\rho_i,\rho_e]$ and $\varphi_0:=\rho_0/\hat r_{G_0}-1$, the amplitude obeys
\begin{equation}\label{eq:realization}
   \norm{\varphi_0}_{L^\infty(S^1)}\ \le\ \frac{\rho_e-\rho_i}{\rho_i}
   \ =\ \frac{\eta}{\rho_i}\ \le\ 2\eta;
\end{equation}
\item\label{it:real-close} $G_0\subset\Etil$ and
\begin{equation}\label{eq:realization-close}
   \abs{\Etil\,\triangle\,G_0}\ =\ \abs{\Etil\setminus G_0}
   \ \le\ C\,\eta\,(\Exc+a\eta).
\end{equation}
\end{enumerate}
Consequently $G_0$ satisfies the hypotheses of \cref{thm:graphG} on the
amplitude, the trapping ball, and the volume normalisation; only the
quadratic mode-one bound \eqref{eq:mode1hyp} remains to be supplied, which is
done at the optimal centre in \cref{sec:main}. No mollification is performed
and no continuity of $\rho_0$ is claimed or used.
\end{proposition}

\begin{proof}
\emph{\ref{it:real-meas}.} The profile $\rho_0=r_1$ is measurable: it is the
essential infimum, along the ray $R_\theta$, of the radii at which the slice
$\Etil_\theta$ first exits $\Etil$, and $\theta\mapsto r_1(\theta)$ is
measurable because for each $t$ the set $\{\theta:r_1(\theta)>t\}=\{\theta:
(0,t)\subset\Etil_\theta\ \text{up to }\HH^1\text{-null}\}$ is, up to a null
set, the slice projection of the open set $\Etil\cap(B_t(z)\setminus\set z)$,
hence measurable (as in the measurability of $R$ in \eqref{eq:Rdef}). By
\eqref{eq:fold-slice} every slice contains the innermost interval $(0,r_1)$
with $r_1\in[\rho_i,\rho_e]$, so $\rho_0\in[\rho_i,\rho_e]$. Passing to the
lower-semicontinuous envelope of $\rho_0$ along rays makes
$G_0=\{r<\rho_0\}$ open; this changes $\rho_0$, hence $\partial G_0$, only on
an $\HH^1$-null set of directions, so $\abs{G_0}$ and all symmetric
differences are unchanged. Then $B_{\rho_i}(z)\subset G_0$ because every ray
contains $(0,\rho_i)\subset(0,r_1)$, and $G_0\subset B_{\rho_e}(z)$ because
$\rho_0\le\rho_e$; this is \eqref{eq:real-trap}, and star-shapedness about
$z$ is immediate from the radial definition.

\emph{\ref{it:real-amp}.} From \eqref{eq:real-trap},
$\pi\rho_i^2\le\abs{G_0}\le\pi\rho_e^2$, so $\hat r_{G_0}=
\sqrt{\abs{G_0}/\pi}\in[\rho_i,\rho_e]$. For every $\theta$ both
$\rho_0(\theta)$ and $\hat r_{G_0}$ lie in $[\rho_i,\rho_e]$, whence
$\abs{\rho_0-\hat r_{G_0}}\le\rho_e-\rho_i=\eta$ and
$\norm{\varphi_0}_\infty=\norm{\rho_0-\hat r_{G_0}}_\infty/\hat r_{G_0}\le
\eta/\rho_i\le2\eta$, using $\rho_i\ge\tfrac12$ (\cref{rem:posterior}).

\emph{\ref{it:real-close}.} By construction the slice of $G_0$ in direction
$\theta$ is $(0,\rho_0(\theta))=(0,r_1(\theta))\subset\Etil_\theta$, so
$G_0\subset\Etil$ up to $\HH^1$-null directions, and
$\abs{\Etil\triangle G_0}=\abs{\Etil\setminus G_0}$. The bound
\eqref{eq:realization-close} is exactly \cref{cor:gd}: the slice of
$\Etil\setminus G_0$ is the union of caps
$(r_2,r_3)\cup\cdots\cup(r_{2m},r_{2m+1})$, of total area
$\le\rho_e\,\eta\int_{S^1}m\,d\theta=\tfrac12\rho_e\,\eta\int_{S^1}(M-1)\,
d\theta\le C\,\eta(\Exc+a\eta)$.
\end{proof}

\begin{remark}[Scales, and the role in the reduction]\label{rem:fold-scales}
Specialising to the inset at level $\hat v\sim\tau=\sqrt\dT$ (so
$\diso\le C\dT/\tau\le C\sqrt\dT$, $\Exc\le C\sqrt\dT$, and by \cref{thm:fmp}
$a\le C\sqrt\diso\le C\dT^{1/4}$), the annulus width is $\eta\le
C(\Exc+a)^{1/2}\le Ca^{1/2}\le C\dT^{1/8}$, an \emph{absolute} small
amplitude. Then \cref{cor:gd} gives
\[
   \gd(\Etil)\le C\eta(\Exc+a\eta)
   \le C\bigl(\dT^{1/8}\dT^{1/2}+\dT^{1/4}\dT^{1/4}\bigr)
   =C\bigl(\dT^{5/8}+\dT^{1/2}\bigr)\le C\sqrt\dT,
\]
and \cref{prop:realization} exhibits the measurable core graph $G_0$ of
amplitude $\norm{\varphi_0}_\infty\le2\eta\le C\dT^{1/8}$ with
$\abs{\Etil\triangle G_0}\le C\sqrt\dT$. These are precisely the scales fed
into the cutting estimate (\cref{lem:cutdeficit}) and the assembly
(\cref{thm:reduction}): the graph defect balances the level-carried closeness
$\abs{\Om\triangle\Etil}\sim\sqrt\dT$ at the optimal exponent, and the
measurable-profile form of \cref{thm:graphG} consumes $G_0$ directly, with no
intervening regularisation.
\end{remark}
